% \iffalse meta-comment
%
%% Copyright (C) 2020-2026 by Marcel Krueger
%%
%% This file may be distributed and/or modified under the
%% conditions of the LaTeX Project Public License, either
%% version 1.3c of this license or (at your option) any later
%% version. The latest version of this license is in:
%%
%% http://www.latex-project.org/lppl.txt
%%
%% and version 1.3 or later is part of all distributions of
%% LaTeX version 2005/12/01 or later.
%
%<*batch>
%<*gobble>
\ifx\jobname\relax\let\documentclass\undefined\fi
\ifx\documentclass\undefined
\csname fi\endcsname
%</gobble>
\input docstrip.tex
\keepsilent
\generate{
  \file{luamml.sty}{\from{luamml.dtx}{package,luatex}}
  \file{luamml-pdf.sty}{\from{luamml.dtx}{package,pdftex}}
}
\endbatchfile
%</batch>
%<*gobble>
\fi
\expandafter\ifx\csname @currname\endcsname\empty
\csname fi\endcsname
%</gobble>
%<*driver>
\DocumentMetadata{tagging=on}
\documentclass{l3in2edoc}
\usepackage{luamml}
\usepackage{csquotes,luacolor,array}
\MakeShortVerb{\|}
\RecordChanges
\ProvideDocElement[printtype=\textit{socket},idxtype=socket,idxgroup=Sockets]{Socket}{socketdecl}
\ProvideDocElement[printtype=\textit{plug},idxtype=plug,idxgroup=Plugs]{Plug}{plugdecl}

\begin{document}
\tracingmathml2
\DocInput{luamml.dtx}
\PrintIndex
\PrintChanges
\end{document}
%</driver>
%<*gobble>
\fi
%</gobble>
% \fi
%
% \GetFileInfo{luamml.sty}
% \title{The \pkg{luamml} package%
%   \thanks{This document corresponds to \pkg{luamml}~\fileversion, dated~\filedate.}%
% }
% \author{Marcel Krüger}
%
% \maketitle
%
% \begin{documentation}
% \section{Use case}
% When generating output for the web or tagged output, mathematical content should often be represented as MathML.
% This uses Lua\TeX~callbacks to automatically attempt to convert Lua\TeX~math mode output into MathML.
%
% \section{Usage}
% The \pkg{luamml} package is designed to be used in automated ways by other packages and usually should not be invoked directly by the end user.
%
% Typically you will use it through \pkg{latex-lab-math}. See the \pkg{latex-lab-math} documentation for user instructions.
%
% If you want to use \pkg{luamml} in another context, the \pkg{luamml} package can be included directly, followed by enclosing blocks which should generate files with \cs{luamml_begin_single_file:} and \cs{luamml_end_single_file:}.
% The filename can be set with \cs{luamml_set_filename:n}.
%
% \section{Improving MathML conversion}
% When using constructs which do not automatically get converted in acceptable form, conversion hints can be provided with \cs{luamml_annotate:en}.
% This allows to provide a replacement MathML structure in Lua table form, for example
% \begin{verbatim}
% \luamml_annotate:en {
%   nucleus = true,
%   core = {[0] = 'mi', 'TeX'},
% }{
%   \hbox{\TeX}
% }
% \end{verbatim}
% produces a |<mi>TeX</mi>| element in the output instead of trying to import \TeX~as a mathematical expression.
%
% It it possible to add a structure around the construct, stash that structure
% and then to tell \cs{luamml_annotate:en} to insert it later inside the math.
% For this the keys \texttt{struct} (which takes a label as argument) or \texttt{structnum}
% (which takes a structure number) can be used. For example
% \begin{verbatim}
% $a = b \quad
%  \tagstructbegin{tag=mtext,stash}\tagmcbegin{}
%  \luamml_annotate:en{nucleus=true,structnum=\tag_get:n{struct_num}}
%  {\mbox{some~text~with~\emph{structure}}}
%  \tagmcend\tagstructend
% $
% \end{verbatim}
% Such a construction should check that the flag for structure elements has actually
% been set to avoid orphaned structures if the stashed structure is ignored.
%
% More about the table structure is explained in an appendix.
%
% \section{Features \& Limitations}
% Currently all mathematical expressions which purely contain Unicode encoded math mode material without embedded non-math should get converted successfully.
% Usage with non-Unicode math (\TeX's 8-bit math fonts) is highly experimental and undocumented.
% Any attempt to build complicated structures by embedding arbitrary \TeX\ code in the middle of math mode needs to have a MathML replacement specified.
% We try to automate more cases in the future.
%
% \appendix
% \input{luamml-algorithm}
% \end{documentation}
%
% \begin{implementation}
% \section{Package Implementation}
% \subsection{Initialization}
% \iffalse
%<*package>
% \fi
%    \begin{macrocode}
%<@@=luamml>
%<*luatex>
\ProvidesExplPackage {luamml} {2026-06-20} {0.9.2}
  {Automatically generate presentational MathML from LuaTeX math expressions}
%</luatex>
%<*pdftex>
\ProvidesExplPackage {luamml-pdf} {2026-06-20} {0.9.2}
  {MathML generation for L̶u̶a̶pdfLaTeX}
%</pdftex>
%    \end{macrocode}
%
% \subsection{Initialization}
% These variable have to appear before the Lua module is loaded and will be used to
% communicate information to the callback.
%
% Here \cs{tracingmathml} does not use a expl3 name since it is not intended for
% programming use but only as a debugging helper for the user.
% The other variables are internal, but we provide public interfaces for setting
% them later.
%    \begin{macrocode}
\int_new:N \l__luamml_flag_int
%    \end{macrocode}
% This integer decides if the xml uses pretty indentation if output to a file
% or to the terminal. It is used in luamml-tex.lua with a bitwise operation. 
% Indentation in files is activated if bit 0 is set.
% Indentation in the log is activated if bit 1 is set.
% Bit 2 sets indentation for code added with the \cs{luamml_register_output_hook:N}.
% Bit 3 sets indentation for MathML written as stream to the PDF with 
% \cs{luamml_get_last_mathml_stream:e}.
%    \begin{macrocode}
\int_new:N \l__luamml_pretty_int
%<luatex>\tl_new:N \l__luamml_filename_tl
\tl_new:N \l__luamml_root_tl
\tl_set:Nn \l__luamml_root_tl { mrow }
\tl_new:N \l__luamml_label_tl
%<pdftex>\int_new:N \g__luamml_formula_id_int
%<luatex>\int_new:N \tracingmathml

\int_set:Nn \l__luamml_pretty_int { 1 }
%    \end{macrocode}
%
% Currently we disable flattening in Lua\TeX{} to avoid annotations getting lost.
% Instead we load \texttt{luamml-mathflatten} which reimplements math flattening
% at a later stage.
%    \begin{macrocode}
%<luatex>\tex_mathflattenmode:D = 0 \scan_stop:
%<luatex>\lua_load_module:n { luamml-mathflatten }
%    \end{macrocode}
%
% Now we can load the Lua module which defines the callback.
% Of course until pdf\TeX starts implementing \cs{directlua} this is only
% done in Lua\TeX.
%    \begin{macrocode}
%<luatex>\lua_load_module:n { luamml-tex }
%    \end{macrocode}
%
% \subsection{Hook}
% We also call a hook with arguments at the end of every MathML conversion with the result.
% Currently only implemented in Lua\TeX{} since it immediately provides the output.
% TODO: the latex-lab-math code registers its own function and so effectively kills the
% hook here. This should be sorted \ldots.
%    \begin{macrocode}
%<*luatex>
\hook_new_with_args:nn { luamml / converted } { 1 }

\cs_new_protected:Npn \__luamml_output_hook:n {
  \hook_use:nnw { luamml / converted } { 1 }
}
\__luamml_register_output_hook:N \__luamml_output_hook:n
%</luatex>
%    \end{macrocode}

%
% \subsection{Flags}
% The most important interface is for setting the flag which controls how the
% formulas should be converted.
%
% \begin{macro}[no-user-doc]{\luamml_process:}
%   Consider the current formula to be a complete, free-standing mathematical
%   expression which should be converted to MathML. Additionally, the formula
%   is also saved in the \texttt{start\_math} node as with
%   \cs{luamml_save:}.
%    \begin{macrocode}
\cs_new_protected:Npn \luamml_process: {
  \tl_set:Nn \l__luamml_label_tl {}
  \int_set:Nn \l__luamml_flag_int { 3 }
}
%    \end{macrocode}
% Temporarily for compatibility
%    \begin{macrocode}
\cs_set_eq:NN \luamml_flag_process: \luamml_process:
%    \end{macrocode}
% \end{macro}
%
% \begin{macro}{\__luamml_maybe_structelem:}
% A internal helper which can be added to a tag to preserve the external state
% of the structelem flag.
%    \begin{macrocode}
\cs_new:Npn \__luamml_maybe_structelem: {
  (
    8 * \int_mod:nn {
      \int_div_truncate:nn { \l__luamml_flag_int } {8}
    } {2}
  ) +
}
%    \end{macrocode}
% \end{macro}
%
% \begin{macro}{\__luamml_style_to_num:N}
%    \begin{macrocode}
\cs_new:Npn \__luamml_style_to_num:N #1 {
%<luatex>  32 * #1
%<*pdftex>
  \token_case_meaning:NnF #1 {
    \displaystyle {0}
    \textstyle {32}
    \scriptstyle {64}
    \scriptscriptstyle {96}
  } {
    \Invalid_mathstyle
  }
%</pdftex>
}
%    \end{macrocode}
% \end{macro}
%
%
% \begin{macro}[no-user-doc]{\luamml_save:n,
%               \luamml_save:nN,
%               \luamml_save:nn,
%               \luamml_save:nNn}
%   Convert the current formula but only save it's representation in the math
%   node without emitting it as a complete formula. This is useful when the
%   expression forms part of a bigger formula and will be integrated into it's
%   MathML tables later by special code.
%   It optionally accepts three parameters: A label, one math style command
%   (\cs{displaystyle}, \cs{textstyle}, etc.) which is the implicit math style
%   (so the style which the surrounding code expects this style to have) and a
%   name for the root element (defaults to \texttt{mrow}).
%   If the root element name is \texttt{mrow}, it will get suppressed in some
%   cases.
%    \begin{macrocode}
\cs_new_protected:Npn \luamml_save:n #1 {
  \tl_set:Nn \l__luamml_label_tl {#1}
  \int_set:Nn \l__luamml_flag_int { \__luamml_maybe_structelem: 1 }
}
\cs_new_protected:Npn \luamml_save:nN #1#2 {
  \tl_set:Nn \l__luamml_label_tl {#1}
  \int_set:Nn \l__luamml_flag_int { \__luamml_maybe_structelem: 17 + \__luamml_style_to_num:N #2 }
}
\cs_new_protected:Npn \luamml_save:nn #1 {
  \tl_set:Nn \l__luamml_label_tl {#1}
  \int_set:Nn \l__luamml_flag_int { \__luamml_maybe_structelem: 5 }
  \tl_set:Nn \l__luamml_root_tl
}
\cs_new_protected:Npn \luamml_save:nNn #1#2 {
  \tl_set:Nn \l__luamml_label_tl {#1}
  \int_set:Nn \l__luamml_flag_int { \__luamml_maybe_structelem: 21 + \__luamml_style_to_num:N #2 }
  \tl_set:Nn \l__luamml_root_tl
}
%    \end{macrocode}
% Temporarily for compatibility
%    \begin{macrocode}
\cs_set_eq:NN \luamml_flag_save:n \luamml_save:n
\cs_set_eq:NN \luamml_flag_save:nN \luamml_save:nN
\cs_set_eq:NN \luamml_flag_save:nn \luamml_save:nn
\cs_set_eq:NN \luamml_flag_save:nNn \luamml_save:nNn
%    \end{macrocode}
% \end{macro}
%
% \begin{macro}[no-user-doc]{\luamml_ignore:}
%   Completely ignore the math mode material.
%    \begin{macrocode}
\cs_new_protected:Npn \luamml_ignore: {
  \int_set:Nn \l__luamml_flag_int { 0 }
}
%    \end{macrocode}
% Temporarily for compatibility
%    \begin{macrocode}
\cs_set_eq:NN \luamml_flag_ignore: \luamml_ignore:
%    \end{macrocode}
% \end{macro}
%
% \begin{macro}[no-user-doc]{\luamml_structelem:}
%   Like \cs{luamml_process:}, but additionally adds PDF structure
%   elements. This only works in Lua\TeX\ and requires that the \pkg{tagpdf} package
%   has been loaded \emph{before} \texttt{luamml}.
%    \begin{macrocode}
%<*luatex>
\cs_new_protected:Npn \luamml_structelem: {
  \tl_set:Nn \l__luamml_label_tl {}
  \int_set:Nn \l__luamml_flag_int { 11 }
}
%    \end{macrocode}
% Temporarily for compatibility
%    \begin{macrocode}
\cs_set_eq:NN \luamml_flag_structelem: \luamml_structelem:
%</luatex>
%    \end{macrocode}
% \end{macro}
%
% \begin{macro}[no-user-doc]{\luamml_set_filename:n}
%   Allows to set a filename to which the generated MathML gets written.
%   Previous content from the file will get overwritten. This includes results
%   written by a previous formula. Therefore this has to be called separately
%   for every formula or it must expand to different values to be useful.
%   The value is fully expanded when the file is written.
%
%   Only complete formulas get written into files (so formulas where
%   \cs{luamml_process:} or \cs{luamml_structelem:} are in effect).
%
%   Only implemented in Lua\TeX, in pdf\TeX\ the arguments for \texttt{pdfmml}
%   determine the output location.
%    \begin{macrocode}
%<*luatex>
\cs_new_protected:Npn \luamml_set_filename:n {
  \tl_set:Nn \l__luamml_filename_tl
}
%</luatex>
%    \end{macrocode}
% \end{macro}
%
% \begin{macro}[no-user-doc]{\luamml_begin_single_file:, \luamml_end_single_file:}
%   Everything between these two commands gets written into the same XML file.
%   The filename is expanded when \cs{luamml_begin_single_file:} gets executed.
%
%   (Implemented in Lua)
% \end{macro}
%
% By default, the flag is set to assume complete formulas.
%    \begin{macrocode}
\luamml_process:
%    \end{macrocode}
%
% \subsection{Annotations}
% These are implemented very differently depending on the engine, but the interface
% should be the same.
% \subsubsection{Lua\TeX}
%    \begin{macrocode}
%<*luatex>
%    \end{macrocode}
% \begin{macro}[no-user-doc]{\luamml_annotate:nen, \luamml_annotate:en}
% A simple annotation scheme: The first argument is the number of top level
% noads to be annotated, the second parameter the annotation and the third
% parameter the actual list of math tokens. The first argument can be omitted to
% let Lua\TeX\ determine the number itself.
%
% Passing the first parameter explicitly is useful for any annotations which
% should be compatible with future pdf\TeX\ versions of this functionality.
%    \begin{macrocode}
\cs_new_protected:Npn \luamml_annotate:nen #1#2#3 {
  \__luamml_flattened_mathgroup_begin:
    #3
  \__luamml_annotate_end:we \tex_numexpr:D #1 \scan_stop: {#2}
}

\cs_new_protected:Npn \luamml_annotate:en #1#2 {
  \__luamml_flattened_mathgroup_begin:
    #2
  \__luamml_annotate_end:e {#1}
}
%    \end{macrocode}
% \end{macro}
%
% \begin{macro}[no-user-doc]{\luamml_attribute:nnn, \luamml_attribute_core:nnn}
% Set a MathML attribute on a MathML node.
%    \begin{macrocode}
\cs_new_protected:Npn \luamml_attribute:een #1#2#3 {
  \__luamml_flattened_mathgroup_begin:
    #3
  \__luamml_annotate_attribute_end:wee {#1} {#2}
}

\cs_new_protected:Npn \luamml_attribute_core:een #1#2#3 {
  \__luamml_flattened_mathgroup_begin:
    #3
  \__luamml_annotate_attribute_end:wee core {#1} {#2}
}
%    \end{macrocode}
% \end{macro}
%
%    \begin{macrocode}
%</luatex>
%    \end{macrocode}

% \subsubsection{pdf\TeX}
%    \begin{macrocode}
%<*pdftex>
%    \end{macrocode}
% \begin{macro}{\__luamml_pdf_showlists:}
% Here and in many other locations the \pdfTeX{} implementation is based on \cs{showlists},
% so we define a internal wrapper which sets all relevant parameters.
%    \begin{macrocode}
\cs_if_exist:NTF \showstream {
  \iow_new:N \l__luamml_pdf_stream
  \iow_open:Nn \l__luamml_pdf_stream { \jobname .tml }
  \cs_new_protected:Npn \__luamml_pdf_showlists: {
    \group_begin:
      \int_set:Nn \tex_showboxdepth:D { \c_max_int }
      \int_set:Nn \tex_showboxbreadth:D { \c_max_int }
      \showstream = \l__luamml_pdf_stream
      \tex_showlists:D
    \group_end:
  }
} {
  \cs_set_eq:NN \l__luamml_pdf_stream \c_log_iow
  \cs_set_eq:NN \__luamml_pdf_set_showstream: \scan_stop:
  \cs_new_protected:Npn \__luamml_pdf_showlists: {
    \group_begin:
      \int_set:Nn \l_tmpa_int { \tex_interactionmode:D }
      \int_set:Nn \tex_interactionmode:D { 0 }
      \int_set:Nn \tex_showboxdepth:D { \c_max_int }
      \int_set:Nn \tex_showboxbreadth:D { \c_max_int }
      \tex_showlists:D
      \int_set:Nn \tex_interactionmode:D { \l_tmpa_int }
    \group_end:
  }
}
%    \end{macrocode}
% \end{macro}
%
%
% \begin{macro}[no-user-doc]{\luamml_annotate:nen, \luamml_annotate:en}
% Now we can define the annotation commands for pdf\TeX.
%    \begin{macrocode}
\cs_generate_variant:Nn \tl_to_str:n { e }
\int_new:N \g__luamml_annotation_id_int
\cs_new_protected:Npn \luamml_annotate:nen #1#2#3 {
  \int_gincr:N \g__luamml_annotation_id_int
  \iow_shipout_x:Nx \l__luamml_pdf_stream {
    LUAMML_MARK_REF:
    \int_use:N \g__luamml_annotation_id_int
    :
  }
  \iow_now:Nx \l__luamml_pdf_stream {
    LUAMML_MARK:
    \int_use:N \g__luamml_annotation_id_int
    :
    count = \int_eval:n {#1},
    #2
    \iow_newline:
    LUAMML_MARK_END
  }
  #3
}
\cs_new_protected:Npn \luamml_annotate:en #1#2 {
  \int_gincr:N \g__luamml_annotation_id_int
  \iow_shipout_x:Nx \l__luamml_pdf_stream {
    LUAMML_MARK_REF:
    \int_use:N \g__luamml_annotation_id_int
    :
  }
  \iow_now:Nx \l__luamml_pdf_stream {
    LUAMML_MARK:
    \int_use:N \g__luamml_annotation_id_int
    :
    count = data.count[\int_use:N \g__luamml_annotation_id_int],
    #1
    \iow_newline:
    LUAMML_MARK_END
  }
  \use:x {
    \iow_now:Nn \l__luamml_pdf_stream {
      LUAMML_COUNT:
      \int_use:N \g__luamml_annotation_id_int
    }
    \__luamml_pdf_showlists:
    \exp_not:n {#2}
    \iow_now:Nn \l__luamml_pdf_stream {
      LUAMML_COUNT_END:
      \int_use:N \g__luamml_annotation_id_int
    }
    \__luamml_pdf_showlists:
  }
}
%    \end{macrocode}
% \end{macro}
%
%    \begin{macrocode}
%</pdftex>
%    \end{macrocode}
%
% \subsection{Trigger for specific formula}
% This only applies for pdf\TeX\ since in Lua\TeX\ everything is controlled by the callback,
% but for compatibility the function is defined anyway.
%
% \begin{macro}[no-user-doc]{\luamml_pdf_write:}
% We could accept parameters for the flag and tag here, but for compatibility
% with Lua\TeX\ they are passed in macros instead.
%    \begin{macrocode}
%<*pdftex>
\cs_new_protected:Npn \luamml_pdf_write: {
  \int_gincr:N \g__luamml_formula_id_int
  \iow_now:Nx \l__luamml_pdf_stream {
    LUAMML_FORMULA_BEGIN:
    \int_use:N \g__luamml_formula_id_int
    :
    \int_use:N \l__luamml_flag_int
    :
    \l__luamml_root_tl
    :
    \l__luamml_label_tl
  }
  \__luamml_pdf_showlists:
  \iow_now:Nx \l__luamml_pdf_stream {
    LUAMML_FORMULA_END
  }
}
%</pdftex>
%<luatex>\cs_new_eq:NN \luamml_pdf_write: \scan_stop:
%    \end{macrocode}
% \end{macro}
%
%    \begin{macrocode}
%    \end{macrocode}
%
% \subsection{Further helpers}
%
% \begin{macro}[no-user-doc]{\RegisterFamilyMapping}
% The Lua version of this is defined in the Lua module.
%    \begin{macrocode}
%<*pdftex>
\NewDocumentCommand \RegisterFamilyMapping {m m} {
  \iow_now:Nx \l__luamml_pdf_stream {
    LUAMML_INSTRUCTION:REGISTER_MAPPING: \int_use:N #1 : #2
  }
}
%</pdftex>
%    \end{macrocode}
% \end{macro}
%
% \subsection{Sockets}
% In various places luamml has to add code to kernel commands. This is done through
% tagging sockets which are predeclared in lttagging.
%
% \begin{tabular}{>{\ttfamily}ll}
% name & number of argument\\\hline
% math/luamml/save/nNn & 1\\
% math/luamml/save/nn & 1\\
% math/luamml/annotate/false & 2\\
% math/luamml/array/save & 0\\
% math/luamml/array/finalize &0\\
% math/luamml/array/initcol & 0\\
% math/luamml/array/finalizecol & 1\\
% math/luamml/mtable/finalizecol & 1\\
% math/luamml/mtable/innertable/save&0\\
% math/luamml/mtable/smallmatrix/save&0\\
% math/luamml/mtable/innertable/finalize&0\\
% math/luamml/mtable/finalize&1\\
% math/luamml/mtable/aligncol&1\\
% math/luamml/mtable/tag/save&0\\
% math/luamml/mtable/tag/set&0\\
% math/display/tag/begin&0\\
% math/display/tag/end&0\\
% math/luamml/hbox&2\\
% math/luamml/artifact&0\\
% math/luamml/finph@nt&2\\   
% math/luamml/finsm@sh&2\\
% \end{tabular}
%
% \subsubsection{Save sockets}
% These sockets are wrappers around the \cs{luamml_save:...} commands
%
%    \begin{macrocode}
\NewTaggingSocketPlug{math/luamml/save/nNn}{luamml}
 {
   \luamml_save:nNn #1
 }
\AssignTaggingSocketPlug{math/luamml/save/nNn}{luamml}
\NewTaggingSocketPlug{math/luamml/save/nn}{luamml}
 {
   \luamml_save:nn #1
 }
\AssignTaggingSocketPlug{math/luamml/save/nn}{luamml}
%    \end{macrocode}
%
% \subsubsection{sockets to annotate content}
%
%    \begin{macrocode}
%<*luatex>
\NewTaggingSocketPlug{math/luamml/annotate/false}{luamml}
 {
   \luamml_annotate:en { core = false } 
    {
      #2
    }
 }
\AssignTaggingSocketPlug{math/luamml/annotate/false}{luamml} 
%</luatex>   
%    \end{macrocode}
% \subsubsection{socket plugs for the array package}
% 
% The luamml support makes only sense with luatex.
%    \begin{macrocode}
%<*luatex>
\AddToHook{package/array/after}{\lua_now:n { require'luamml-array' }}
%    \end{macrocode}
% \begin{plugdecl}{tagsupport/math/luamml/array/save}
%  The socket of this plug is used in \cs{endarray}.
%    \begin{macrocode}
\NewTaggingSocketPlug{math/luamml/array/save}{luamml}
 {
   \__luamml_array_save_array:
 }
%    \end{macrocode}
% \end{plugdecl}
%
% \begin{plugdecl}{tagsupport/math/luamml/array/finalize}
%  This socket of this plug is used in \cs{endarray}.
%    \begin{macrocode}
\NewTaggingSocketPlug{math/luamml/array/finalize}{luamml}
 {
   \mode_if_math:T { \__luamml_array_finalize_array: }
 }
%    \end{macrocode}
% \end{plugdecl}
%
% \begin{plugdecl}{tagsupport/math/luamml/array/initcol}
%  The socket of this plug is used in \cs{@classz}.
%    \begin{macrocode}
\NewTaggingSocketPlug{math/luamml/array/initcol}{luamml}
 {
   \__luamml_array_init_col:
 }
%    \end{macrocode}
% \end{plugdecl}
%
%
% \begin{plugdecl}{tagsupport/math/luamml/array/finalizecol}
%  The socket of this plug is used in \cs{@classz}.
%    \begin{macrocode}
\NewTaggingSocketPlug{math/luamml/array/finalizecol}{luamml}
 {
   \__luamml_array_finalize_col:w #1~
 }
%    \end{macrocode}
% \end{plugdecl}
%    \begin{macrocode}
\AssignTaggingSocketPlug{math/luamml/array/save}{luamml}
\AssignTaggingSocketPlug{math/luamml/array/finalize}{luamml}
\AssignTaggingSocketPlug{math/luamml/array/initcol}{luamml}
\AssignTaggingSocketPlug{math/luamml/array/finalizecol}{luamml}
%</luatex>
%    \end{macrocode}
% \subsubsection{amsmath alignments}
%
% This socket is used at the end of alignment cells and adds the content to
% the current row.
%
%    \begin{macrocode}
%<*luatex>
\NewTaggingSocketPlug{math/luamml/mtable/finalizecol}{luamml}
 {
   \use:c{__luamml_amsmath_add_#1_to_row:}
 }
\AssignTaggingSocketPlug{math/luamml/mtable/finalizecol}{luamml}

%</luatex>
%    \end{macrocode}
%
% These sockets save an inner table
%
%    \begin{macrocode}
%<*luatex>
\NewTaggingSocketPlug{math/luamml/mtable/innertable/save}{luamml}
 {
   \__luamml_amsmath_save_inner_table:n \@currenvir
 }
\AssignTaggingSocketPlug{math/luamml/mtable/innertable/save}{luamml}
\NewTaggingSocketPlug{math/luamml/mtable/smallmatrix/save}{luamml}
 {
   \__luamml_amsmath_save_smallmatrix:
 }
\AssignTaggingSocketPlug{math/luamml/mtable/smallmatrix/save}{luamml}
\NewTaggingSocketPlug{math/luamml/mtable/innertable/finalize}{luamml}
 {
   \__luamml_amsmath_finalize_inner_table:
 }
\AssignTaggingSocketPlug{math/luamml/mtable/innertable/finalize}{luamml}
%</luatex>
%    \end{macrocode}
%
%
% This socket finalize the \texttt{mtable} in alignments like align or gather.
% It takes an argument, the environment.
% It should be used normally with \cs{UseExpandableTaggingSocket}.
%
%    \begin{macrocode}
%<*luatex>
\NewTaggingSocketPlug{math/luamml/mtable/finalize}{luamml}
 {
   \__luamml_amsmath_finalize_table:n {#1}
 }
\AssignTaggingSocketPlug{math/luamml/mtable/finalize}{luamml}
%</luatex>
%    \end{macrocode}
%
% This socket adds attributes for the alignment in \texttt{multline}.
% It takes an argument, the alignment.
%
%    \begin{macrocode}
%<*luatex>
\NewTaggingSocketPlug{math/luamml/mtable/aligncol}{luamml}
 {
   \__luamml_amsmath_set_row_columnalign:n {#1}
 }
\AssignTaggingSocketPlug{math/luamml/mtable/aligncol}{luamml}
%</luatex>
%    \end{macrocode}

%
% \subsubsection{Tags and labels}
% These sockets save and set tags and labels in alignments.
%
%    \begin{macrocode}
%<*luatex>
\int_new:N \l__luamml_amsmath_tag_struct_num_int
\NewTaggingSocketPlug{math/luamml/mtable/tag/save}{luamml}
 {
   \__luamml_amsmath_save_tag_with_struct_elem:N \l__luamml_amsmath_tag_struct_num_int
 }
\AssignTaggingSocketPlug{math/luamml/mtable/tag/save}{luamml}
\NewTaggingSocketPlug{math/luamml/mtable/tag/set}{luamml}
 {
   \__luamml_amsmath_set_tag:
 }
\AssignTaggingSocketPlug{math/luamml/mtable/tag/set}{luamml}

%</luatex>
%    \end{macrocode}
%
% If math structure elements are created the Lbl-structure of a tag
% must be moved inside the math structure, typically as an additional column in an
% \texttt{mtable} with an intent \texttt{:equation-label} or \texttt{:no-equation-label}.
% 
% The luamml-code handles this by stashing the Lbl-structure, storing the
% structure number in an array and reusing it once it creates the math structure elements.
% 
% This should only be done for specific environments, we define
% a constant to test: 
%
%    \begin{macrocode}
%<*luatex>
\clist_map_inline:nn 
  {
    eqnarray,
    eqnarray*,
    align,
    align*,
    alignat,
    alignat*,
    xalignat,
    xalignat*,
%    \end{macrocode}
%  there is never a tag/label in xxalignat, so does it make sense to add a label column? 
%  Left out for now.
%    \begin{macrocode}
    %xxalignat, 
    flalign,
    flalign*,
    gather,
    gather*,
%    \end{macrocode}
% multline have at most one tag, so we do not use a label column
% but rely on the external Lbl for now.
%    \begin{macrocode}
    %multline, % NO
    %multline*, % NO
    equation,
    equation*,
%    \end{macrocode}
%   split has never a numbering so is ignored
%    \begin{macrocode}
    %split, % NO
  }
  {\tl_const:cn { c__luamml_label_#1_tl}{}}
%    \end{macrocode}
% 
% 
%    \begin{macrocode}
\NewTaggingSocketPlug{math/display/tag/begin}{luamml}
 { 
   \tag_mc_end: 
   \bool_lazy_and:nnT
    { \tl_if_exist_p:c { c__luamml_label_ \@currenvir _tl } }
    { \int_if_odd_p:n { \int_div_truncate:nn { \l__luamml_flag_int } { 8 } } }
    {   
      \tag_struct_begin:n {tag=mtext,stash}
      \int_set:Nn \l__luamml_amsmath_tag_struct_num_int { \tag_get:n {struct_num} }
    }   
   \tag_mc_begin:n {}
 }
\AssignTaggingSocketPlug{math/display/tag/begin}{luamml} 
\NewTaggingSocketPlug{math/display/tag/end}{luamml}
 {
   \tag_mc_end:
   \bool_lazy_and:nnT
    { \tl_if_exist_p:c { c__luamml_label_ \@currenvir _tl } }
    { \int_if_odd_p:n { \int_div_truncate:nn { \l__luamml_flag_int } { 8 } } }
    {
      \tag_struct_end:
    }
   \tag_mc_begin:n {}
 }
\AssignTaggingSocketPlug{math/display/tag/end}{luamml}
%</luatex>
%    \end{macrocode}
%
%
% \subsubsection{Horizontal boxes}
% This socket annotates an \cs{hbox} inside box commands used in math.
% We test for the socket until the release 2025-06-01.
%    \begin{macrocode}
%<*luatex>
\NewTaggingSocketPlug{math/luamml/hbox}{luamml}
 {
   \bool_lazy_and:nnTF
    { \mode_if_math_p: }
    { \int_if_odd_p:n { \int_div_truncate:nn { \l__luamml_flag_int } { 8 } } }
    {
      \tag_struct_begin:n
       {
         tag=mtext,
         stash,
       }
      \tag_mc_begin:n {}
      \luamml_annotate:en
       {
         nucleus = true,
         structnum=\tag_get:n{struct_num}
       }
       { #2 }
      \tag_mc_end:
      \tag_struct_end:
    }
    { #2 }
  }
\AssignTaggingSocketPlug{math/luamml/hbox}{luamml}
%</luatex>
%    \end{macrocode}
%
% \subsubsection{Artifact characters}
% Unicode characters like a root sign should be marked as artifacts
% to avoid duplication e.g. in derivation if mathml
% structure elements are used that imply the meaning.
%    \begin{macrocode}
%<*luatex>
\NewTaggingSocketPlug{math/luamml/artifact}{luamml}
 {
   \int_if_odd:nT { \int_div_truncate:nn { \l__luamml_flag_int } { 8 } }
    {
      \tag_mc_begin:n{artifact}
    }
 }
\AssignTaggingSocketPlug{math/luamml/artifact}{luamml}
%</luatex>
%    \end{macrocode}
%
% \subsubsection{Math phantom socket}
% This socket is used around \cs{finph@nt}.
%
%    \begin{macrocode}
%<*luatex>
\NewTaggingSocketPlug{math/luamml/finph@nt}{luamml}
 {
   \luamml_annotate:nen {1}
    {
     nucleus = true,
     mathml =
      {
        [0] = 'mpadded',
        \ifh@\else
         width = 0,
        \fi
        \ifv@\else
         height = 0, depth = 0,
        \fi
        consume_label'mathphant',
      },
    }
    { #2 }
 }
\AssignTaggingSocketPlug{math/luamml/finph@nt}{luamml}
%</luatex>
%    \end{macrocode}

% \subsubsection{Math smash socket}
% This socket is used around \cs{finsm@sh}.
% 
%    \begin{macrocode}
%<*luatex>
\NewTaggingSocketPlug{math/luamml/finsm@sh}{luamml}
 {
   \luamml_annotate:nen {2}
    {
     nucleus = true,
     mathml = consume_label('mathsmash',
           function(padded, core)
             padded.height, padded.depth = 0, 0~
             return~padded, core~
           end),
    }
    { #2 }
 }
\AssignTaggingSocketPlug{math/luamml/finsm@sh}{luamml}
%</luatex>
%    \end{macrocode}
% \subsection{Patching}
% For some packages, we ship with patches to make them more compatible and to
% demonstrate how other code can be patched to work with \texttt{luamml}.
%
% These are either loaded directly if the packages are loaded or delayed using
% \LaTeX's hook system otherwise.
% \begin{macro}{\__luamml_patch_package:nn, \__luamml_patch_package:n}
% For this, we use two helpers: First a wrapper which runs arbitrary code either
% now (if the package is already loaded) or as soon as the package loads, second
% an application of the first one to load packages following \texttt{luamml}'s
% naming scheme for these patch packages.
%    \begin{macrocode}
\cs_new_protected:Npn \__luamml_patch_package:nn #1 #2 {
  \@ifpackageloaded {#1} {#2} {
    \hook_gput_code:nnn {package/#1/after} {luamml} {#2}
  }
}
\cs_new_protected:Npn \__luamml_patch_package:n #1 {
  \__luamml_patch_package:nn {#1} {
    \RequirePackage { luamml-patches-#1 }
  }
}
%    \end{macrocode}
% \end{macro}
%
% We currently provide minimal patching for the kernel, \pkg{amsmath}.
% Currently only the kernel code supports pdf\TeX, but it's planned to extend this.
%    \begin{macrocode}
\RequirePackage { luamml-patches-kernel }
%<*luatex>
\__luamml_patch_package:n {amsmath}
%</luatex>
%    \end{macrocode}

% \iffalse
%</package>
% \fi
% \end{implementation}
% \Finale
